\section{Customer-Owned Direct Remote Route}
\label{sec:pilot-mobile-direct-remote-route}

Pilot now provides a direct remote phone-to-mother route as the coded completion of
\texttt{MOB-0307}.  The route does not copy the mother brain, Personal data, Inbox,
Workspace or Boardroom state into a Tessaris service.  It allows the same signed mobile
operations used on the trusted home network to reach a customer-controlled HTTPS
endpoint when the local endpoint is genuinely unreachable.

\subsection{Configuration and pairing}

An owner may configure up to four externally reachable HTTPS endpoints through the
mother environment setting \texttt{PILOT\_REMOTE\_ENDPOINTS}.  Pilot rejects HTTP,
embedded credentials, paths, query strings and fragments.  The local endpoint is
removed from the remote set and duplicates are suppressed.

During successful local pairing, the mother issues a bounded descriptor containing the
mother identifier, public-key fingerprint, ordered endpoint set, transport and
application-authentication requirements, issue and expiry times, and a random nonce.
The descriptor states that no relay is used and no mother secret is present.  Its
canonical payload hash and Ed25519 mother signature bind every field.  The descriptor
is delivered only with the successful paired-phone response and stored with that
phone's protected connection record.

The route is configuration, not an automatic router modification.  The owner or
organization remains responsible for a valid public certificate, DNS and an authorized
customer-controlled reverse proxy or direct service exposure.  Automatic rendezvous,
relay fallback and Wi-Fi/mobile-data roaming are separate later tasks.

\subsection{Connection selection}

For every ordinary mobile request, the phone follows this order:
\begin{enumerate}
  \item call the paired trusted-LAN endpoint;
  \item if and only if the network connection itself fails, inspect the cached remote
        descriptor;
  \item verify descriptor expiry, schema, endpoint bounds, mother identity, payload
        signature, no-relay declaration and absence of embedded secrets;
  \item connect to a signed candidate over browser-validated HTTPS;
  \item request that endpoint's fresh signed mother descriptor;
  \item require the original mother identifier and public-key fingerprint;
  \item submit the unchanged signed mobile operation.
\end{enumerate}

An authorization error, validation error or application failure on the local mother is
not treated as a reason to try another endpoint.  This prevents failover from bypassing
policy and avoids duplicating a consequential request merely because the mother refused
it.  Existing idempotency keys continue to protect operations if a genuine network
failure leaves the client uncertain.

Let $D$ be the signed direct-route descriptor, $M$ the mother identity established at
pairing, $E$ a candidate endpoint and $Q$ the already possession-signed request.  Remote
submission is allowed only when
\[
  \mathrm{Signature}_{M}(D) \land \mathrm{Fresh}(D) \land E\in D
  \land \mathrm{TLS}(E) \land \mathrm{Identity}(E)=M \land \mathrm{Valid}(Q).
\]
The network endpoint therefore transports authority but cannot create it.

\subsection{Security and ownership properties}

TLS 1.2 or newer protects the direct channel, while the existing phone certificate,
short-lived capability lease and Ed25519 request signature provide application-level
mutual authentication.  A DNS, QR or endpoint substitution cannot succeed without both
the original mother's signature and its private identity key.  Revoking the phone or
lease continues to reject requests on local and remote addresses alike because both
addresses terminate at the same authoritative mother.

The descriptor contains no credentials, provider tokens or mother-private key.  It is
valid for at most 72 hours and normally 24 hours.  The route uses no central Pilot
plaintext store and declares \texttt{relay\_used=false}.  A future unreadable relay may
be selected only through the separately governed \texttt{MOB-0308} design.

\subsection{Verification and remaining work}

Automated tests cover signed route creation, unsafe endpoint rejection, bounded endpoint
sets, secret absence, descriptor tampering, pairing delivery, LAN-first ordering, mother
substitution refusal and routing of existing mobile requests through the failover
boundary.  All 560 AION Fabric tests pass.  The optimized phone build type-checks and
generates the mobile route, and this append-only section compiles cleanly.

This is coded direct-route completion, not a claim that the current household router or
public Internet endpoint has been configured.  \texttt{MOB-0308} remains responsible for
an end-to-end encrypted unreadable relay fallback, and \texttt{MOB-0309} remains
responsible for automatic Wi-Fi/mobile-data roaming, bounded reconnect and field proof.
